\scanpage{108}
% archive sheet number: 107
% source 01-04, top
\uline{NB} Est-il clair que la condition $\alpha^2+\alpha-1=0$
est \uline{\unsure{suffisante|suffisant}} pour avoir une représentation
de l'icosaèdre ? On aura une représentation affine du gr.
% source 05-06, top
\[
\Gamma'=\left\{\sigma_0,\sigma_1,\sigma_2\ \left|\
\begin{aligned}
\sigma_0^2=\sigma_1^2=\sigma_2^2&=(\sigma_0\sigma_1)^3=(\sigma_1\sigma_2)^{\unsure{6|5}}\\
&=(\sigma_0\sigma_2)^2=1
\end{aligned}\right.\right\}.
\]
% source 07-09, upper middle
\unsure{Ce|Rque} groupe \deleted{\unsure{présenté}} admet $\Gamma$ comme
groupe quotient. Si on \deleted{\unsure{savait}}\inserted{utilise}
$\Gamma'\simeq\Gamma$ \unsure{---} \unsure{qu'on} connaît.
% source 10-13, middle
Sinon il faut une vérification un peu emmerdante
(mais de pure routine) que \deleted{\unsure{les}}\inserted{l'ens.}
des six points $s_0,s_1,s_2,s_3$, $s'_3,s''_3$ et leurs
``antipodiques'' \unsure{définis|definis} ``a priori'' \unsure{sont|est} stable sous
$\sigma_0,\sigma_1,\sigma_2$ (i.e. sous $\Gamma'$) et
\deleted{\unsure{que}} plus précisément que les op. des
$\sigma_0,\sigma_1,\sigma_2$ sur l'ens. de ces $12$ pts ``est'' la
``$=$'' que l'op. sur les $12$ pts de l'icos. comb. $\unsure{\Pi|T}$ ....
% source 14, separator
\figure{\begin{tikzpicture}\draw (0,0)--(2,0);\end{tikzpicture}}
% source 15-16, bottom
Ici $\sigma_0,\sigma_1,\sigma_2$ sont des réflexions. Les
hyperplans correspondants sont \deleted{\unsure{en}}
% source 17-19, bottom
\[
\left\{\begin{aligned}
H_0&:\quad 2x+y+z-1=0\\
H_1&:\quad -x+y+\alpha'z=0\\
H_2&:\quad \deleted{\unsure{\alpha|a}}y-z=0
\end{aligned}\right.
\]
